% ID: SEL-D0018-BC-E05
% ESTADO: cotejada. Segunda pasada visual completa y borrador compilado revisados.
% Fuente: Selectividad/matematicas_aplicadas/matematicas_aplicadas_P_F.pdf
% SHA-256: fe6fd5b48407569d41169c80d71e6997050d6d06d1afd41cd639a553f9b712b5
% Página: 2; bloque C; ejercicio original 5.

Se sabe que el 65\% de los estudiantes de bachillerato de Andalucía ha participado en programas Erasmus+ y que de ellos, el 80\% ha mejorado su calificación en lengua extranjera. De los estudiantes que no han participado en programas Erasmus+, mejoran su calificación en lengua extranjera el 30\%. Se elige al azar un estudiante de bachillerato de Andalucía.
\begin{enumerate}[label=\alph*)]
\item \textbf{(1.5 puntos)} \textquestiondown{}Cuál es la probabilidad de que haya mejorado su calificación en lengua extranjera?
\item \textbf{(1 punto)} Si se sabe que ha mejorado su calificación en lengua extranjera, \textquestiondown{}cuál es la probabilidad de que haya participado en un programa Erasmus+?
\end{enumerate}
